\documentclass[12pt, twoside]{article}
\input{/Users/chris/github/course-files/docs/ib/preamble}

\fancyhead[LE]{\thepage}
\fancyhead[RO]{Name: \hspace{4cm} \,\\ \hspace{2.9cm} \,}
\fancyhead[LO]{La Scuola d'Italia / Dr.\ Huson / IB Math HL \\* 28 May 2026}
\fancyfoot[C]{\thepage}

\begin{document}
\subsubsection*{8.1 Classwork: Rational Functions (HL)}

\noindent\emph{A graphing calculator is allowed. Answer on lined or graph paper;
label each part clearly. Show enough algebra that a marker can follow your reasoning
--- a GDC answer alone is not sufficient unless the part says ``write down''.}

\begin{enumerate}[itemsep=0.25cm]

%--- Problem 1: Quadratic over linear, oblique asymptote [12 marks] ---
%--- Modeled after AA HL 2023 Nov P2 TZ1 Q11 (qb_id 23N.2.AHL.TZ1.11) ---
\item[1.] [Maximum mark: 12] \ The function $f$ is defined by
$\displaystyle f(x)=\dfrac{x^{2}-3x-4}{x-2}, \quad x\in\mathbb{R},\ x\ne 2.$

\begin{enumerate}[itemsep=0.1cm]
  \item Write down the equation of the vertical asymptote of the graph of $f$.
  \hfill [1]
  \item Find the coordinates of the points where the graph of $f$ crosses the
  $x$-axis, and the $y$-axis. \hfill [3]
  \item Using polynomial division, or otherwise, show that
  $\displaystyle f(x)=x-1-\dfrac{6}{x-2}$, and hence state the equation of the
  oblique asymptote of $f$. \hfill [4]
  \item Using your GDC, sketch the graph of $y=f(x)$ for $-4\le x\le 8$. Label
  the asymptotes (dashed) and the three intercepts. \hfill [2]
  \item Solve the inequality $f(x)>0$. \hfill [2]
\end{enumerate}

\medskip

%--- Problem 2: Reciprocal of a quadratic [10 marks] ---
%--- Modeled after AA HL 2022 May P1 TZ2 Q11 (qb_id 22M.1.AHL.TZ2.11) ---
\item[2.] [Maximum mark: 10] \ Let $\displaystyle f(x)=\dfrac{1}{x^{2}-2x-8}$,
\ $x\in\mathbb{R}$, $x\ne -2$, $x\ne 4$.

\begin{enumerate}[itemsep=0.1cm]
  \item State the equations of the two vertical asymptotes and the horizontal
  asymptote of the graph of $f$. \hfill [3]
  \item Find the $y$-intercept of $f$. \hfill [1]
  \item Show that $f$ has a single local extremum, and find its coordinates.
  State whether it is a local maximum or a local minimum, and justify your
  answer. \hfill [4]
  \item Sketch the graph of $y=f(x)$, indicating asymptotes (dashed), the
  $y$-intercept, and the local extremum. \hfill [2]
\end{enumerate}

\newpage

%--- Problem 3: Linear over quadratic, sign analysis [8 marks] ---
%--- Modeled after AA HL 2021 May P2 TZ1 Q11 (qb_id 21M.2.AHL.TZ1.11) ---
\item[3.] [Maximum mark: 8] \ The function $f$ is defined by
$\displaystyle f(x)=\dfrac{x-2}{x^{2}-9}, \quad x\in\mathbb{R},\ x\ne\pm 3.$

\begin{enumerate}[itemsep=0.1cm]
  \item State the equations of all three asymptotes of the graph of $f$. \hfill [3]
  \item Find the coordinates of the $x$-intercept and the $y$-intercept of $f$.
  \hfill [2]
  \item Solve the inequality $f(x)\ge 0$. Express your answer using interval
  notation. \hfill [3]
\end{enumerate}

\medskip

%--- Problem 4: Inverse of a Mobius + asymptote duality [6 marks] ---
%--- Modeled after AA HL 2021 Nov P1 TZ0 Q02 (qb_id 21N.1.AHL.TZ0.2) ---
\item[4.] [Maximum mark: 6] \ The function $g$ is defined by
$\displaystyle g(x)=\dfrac{ax+4}{3-x}, \quad x\in\mathbb{R},\ x\ne 3,\ a\in\mathbb{R}.$

\begin{enumerate}[itemsep=0.1cm]
  \item Given that $g(x)=g^{-1}(x)$ for all $x$ in the domain of $g$, determine
  the value of $a$. \hfill [4]
  \item For this value of $a$, write down the equation of the line of symmetry
  of the graph of $y=g(x)$, and explain why every self-inverse function has its
  graph symmetric about this line. \hfill [2]
\end{enumerate}

\end{enumerate}
\end{document}
